\chapter{波模}

\section{波的反射}

本章将讨论把波动限制在某一有限区域中时所产生的一些值得注意的现象。首先，我们将介绍几个有关振动弦所特有的事实作为例子，然后由这些事实归纳出一条原理，这一原理很可能是影响最广泛的数学物理原理。

约束波的第一个例子是使波动在一个界面上受到限制。现在来看一下弦上的一维波动这个简单例子。我们也完全可以讨论朝着墙传播的一维声波或其他性质类似的情况，但是弦的例子目前对我们来说已足够了。假定弦的一端固定，例如把弦系于一堵“无限坚实”的墙上。从数学上讲，这种情况可以表述成：弦在位置$x=0$处的位移$y$必须是零，因为该端点没有运动。现在，如果对于墙来说不是这样，那么我们知道运动的通解是$F(x -ct)$和$G(x+ct)$这两个函数之和，前面一项表示弦中沿一个方向传播的波，后面一项表示弦中沿相反方向传播的波，即
\begin{equation}
    \label{Eq:I:49:1}
    y=F(x-ct)+G(x+ct)
\end{equation}
是对任意弦的通解。但是，接下来我们必须满足弦的一端不运动这个条件。如果令方程式（\ref{Eq:I:49:1}）中的$x=0$，并对任何$t$的数值考察$y$的值，于是我们得到
\begin{equation*}
    y=F(-ct)+G(+ct)
\end{equation*}
如果对所有的时间，这个值都是零，那就意味着函数$G(ct)$一定等于$-F(-ct)$。换句话说，任何量的函数$G$必定等于该量负值的函数$F$之负值。如果将这个结果代回到方程（\ref{Eq:I:49:1}）中，我们就发现问题的解是
\begin{equation}
    \label{Eq:I:49:2}
    y=F(x-ct)-F(-x-ct)
\end{equation}
如果令$x =0$，就得到$y =0$，这一点很容易验证。

图\ref{figure:把波的反射看作为两个行波的叠加}表示在$x =0$附近沿负$x$方向传播的一个波，以及另一符号相反、在原点另一边沿正$x$方向传播的假想波。我们说假想的，当然是因为没有弦在原点的另一边振动。
\begin{wrapfigure}{r}{0.35\textwidth}
    \centering
    \includegraphics[width=0.33\textwidth]{Chapter49/把波的反射看作为两个行波的叠加}
    \caption{\footnotesize 把波的反射看作为两个行波的叠加}
    \label{figure:把波的反射看作为两个行波的叠加}
\end{wrapfigure}
弦上总的运动被看成为这样两个波在正$x$区域中的叠加。当它们到达原点$x=0$处，这两个波总是相互抵消，最后，在正$x$区域将只存在第二个（被反射的）波，此波当然将沿相反方向传播。这些结果与下列表述是一致的：如果一个波到达弦的固定端，它将被反射，同时改变符号。对这种反射总是可以这样来理解，即想象向弦端点行进的波，在墙后倒转波形再传出来。简单地说，如果我们假定弦是无限长的，每当我们有一个沿某一方向行进的波时，总有另一个沿相反方向行进的波，而且具有上述的对称性，则在$x=0$处的位移总是零，这时，即使在$x=0$处把弦固定也是毫无影响的。

下面讨论周期波的反射。假定用$F(x-ct)$表示的波是一个正弦波，并且已被反射；那么，反射波$-F(-x-ct)$也是同一个频率的正弦波，但沿相反方向传播。这种情况可以很简单地用复变函数记号来描写：$F(x-ct)=e^{i\omega(t-x/c)}$和$F(-x-ct)=e^{i\omega(t+x/c)}$。如果将这两个表示式代入式（\ref{Eq:I:49:2}），并且令$x$等于0，就可以看到，对所有的$t$值，$y=0$，所以它满足所需的条件。由于指数的性质，式（\ref{Eq:I:49:2}）可以写成更为简单的形式
\begin{equation}
    \label{Eq:I:49:3}
    y =e^{i \omega t}(e^{-i \omega x / c }-e^{i \omega x / c })=-2ie^{i \omega t}\sin(\omega x /c )
\end{equation}

这个解告诉我们一些有趣的、新的东西，这就是如果我们在任一固定的$x$处观察弦，那么它以频率$\omega$振动。不论这一点在哪里，频率都相同！但是在有一些地方，特别是在$\sin(\omega x /c )=0$的地方，则根本没有位移。另外，如果在任一时刻$t$，我们给振动弦拍一张快照，那么将得出一张正弦波的像。但是这个正弦波的位移与时间$t$有关。从方程式（\ref{Eq:I:49:3}）可以看出此正弦波一周的长度等于两个叠加波中任何一个波的波长
\begin{equation}
    \label{Eq:I:49:4}
    \lambda = 2\pi c/\omega
\end{equation}
没有运动的那些点满足条件$\sin(\omega x /c)=0$，这就意味着$(\omega x /c )=0,~\pi,~2\pi,~\cdots,~n \pi,~\cdots$这些点称为波节。在任何两个相邻波节之间，每一点按正弦规律上下运动，但是运动的图样在空间是固定不动的。这一特征就是我们称之为波模 的基本特征。如果能够找到一个运动图样，它具有这样的性质，即在任何点处物体完全按正弦规律运动，而且所有的点都以相同频率运动（虽然有些点比另一些点运动得更多），那么我们就得到了所谓的波模。

\section{具有固有频率的约束波}
下一个有趣的问题是讨论如果弦的两端（比如说在$x=0$和$x=L$处）都固定时所发生的情况。我们可以由波的反射的概念着手，首先考察沿某一方向行进的某种“隆起部分”的运动。随着时间的推移，我们可以预期“隆起部分”将到达一个端点附近，而随着时间的继续推移，由于与来自另一边、方向与符号都相反的假想“隆起部分”合成的结果，它将变成一种小的颤动。最后，原始的“隆起部分”将消失，假想的“隆起部分”将沿反方向运动，并在另一端点重复上述过程。这个问题的解很容易，但是一个有趣的问题是我们能否得到一个正弦运动（刚才描述的解是周期性的 ，但当然不是正弦 周期性的）。让我们来试一下在弦上产生一个正弦周期波。如果弦的一端被系住，我们知道它一定与我们先前的解式（\ref{Eq:I:49:3}）相同。如果弦的另一端被系住，那么它在那一端出现的情况也必定相同。因此，对于周期的正弦运动，唯一的可能性是正弦波必须恰好能适合弦的长度。如果正弦波不能适合弦的长度，那么它的频率就不是一个可以使弦连续不断振动的固有频率。简单地说，如果弦恰好以合适的正弦波形状开始运动，那么它将继续保持完善的正弦波形，并在某个频率上谐和地振动。

数学上，我们可以将上述波形写成$\sin kx$，式中$k$等于方程式（\ref{Eq:I:49:3}）和（\ref{Eq:I:49:4}）中的因子$\omega /c$，这一函数在$x=0$处等于零。但是，它在另一端点也必须等于零。这一点的意义是，这里$k$不再像在一端固定的弦中那样是任意的了。当弦的两个端点均固定时，唯一的可能性就是$\sin(kL)=0$，因为这是使两个端点均保持固定的唯一条件。现在，为了要使正弦值等于零，角度必须等于0，$\pi$，$2\pi$或者另一些$\pi$的整数倍。因此，方程
\begin{equation}
    \label{Eq:I:49:5}
    kL = n\pi
\end{equation}
将给出任何一个可能的$k$值，取决于所取的整数值是多少。对每一个$k$值，就有一个确定的频率$\omega$，按照（\ref{Eq:I:49:3}）它就是
\begin{equation}
    \label{Eq:I:49:6}
    \omega = kc = n\pi c/L
\end{equation}

因此，我们得出下述结论：弦具有一种可以作正弦运动的性质，但仅能以某些确定的频率作正弦运动 。这是约束波的最重要的特征。无论系统怎么复杂，结果表明总是存在某些运动图样，它们具有完好的对时间的正弦依赖关系，但频率则取决于该特定系统及其边界的性质。对弦来说，我们可以有许多不同的频率，根据定义，每一个频率对应一个模式，因为模式就是按正弦方式反复重演的运动图样。图\ref{figure:振动弦的前三种模式}表示弦的前三种模式。对于第一种模式，波长$\lambda$等于$2L$。只要将波延长到$x=2L$，得出一个完整的正弦波，就能看出这一点。通常，角频率$\omega$等于$2\pi c$除以波长，现在由于$\lambda$等于$2L$，角频率就等于$\pi c /L$，与式（\ref{Eq:I:49:6}）中$n=1$的结果相符。我们把第一种模式的频率称为$\omega_1$。第二种模式给出了两个波腹，在其中点有一个波节。于是，这种模式的波长就等于$L$。相应的$k$值增大为两倍，频率也增大为两倍，等于$2\omega_1$。对于第三种模式，频率等于$3\omega_1$，如此等等。因此弦的所有不同频率都是最低频率$\omega_1$的整数倍——1倍、2倍、3倍、4倍，等等。

\begin{figure}[htbp]
    \centering
    \begin{minipage}[t]{0.42\textwidth}
        \centering
        \includegraphics[width=6.2cm]{Chapter49/振动弦的前三种模式}
        \caption{振动弦的前三种模式}
        \label{figure:振动弦的前三种模式}
    \end{minipage}
    \quad\quad
    \begin{minipage}[t]{0.42\textwidth}
        \centering
        \includegraphics[width=6.2cm]{Chapter49/两种模式组合成一个行波}    
        \caption{两种模式组合成一个行波}
        \label{figure:两种模式组合成一个行波}
    \end{minipage}
\end{figure}

现在，回到弦的一般运动，结果是任何可能的运动总可以这样来分析，即认为弦上同时振动着多个模式。事实上，对一般运动来说，必须同时激发无数个模式。为了得到有关这方面的某些概念，让我们用图来说明，当有两个模式同时振动时所发生的情况：假定第一种模式按图\ref{figure:两种模式组合成一个行波}中所示的次序振动着，图中画出在最低频率的半个周期内每隔相等的时间间隔弦的挠曲形状。

现在，我们假定同时还存在着两种模式的振动。图\ref{figure:两种模式组合成一个行波}也画出了一系列这种模式的图样，该模式在开始时与第一种模式有90$^\circ$的相位差。这意味着它开始时没有位移，但是弦的两半具有相反方向的速度。现在，我们回想一下有关线性系统的一般原理：如果存在任意两个解，那么它们的和也是该系统的一个解。因此，由图\ref{figure:两种模式组合成一个行波}所示的两个解相加得到的位移应是弦的第三种可能的运动。合成结果也画在图中，它开始显示出在弦的两个端点之间来回运动的一个“隆起部分”的样子，虽然只有两个模式，我们作不出非常好的图，好的图需要更多的模式。事实上，这个结果是关于线性系统的重要原理的一个特殊情况：

任何运动都可以这样来分析，即设想它是所有各种由适当振幅和相位组成的不同模式的运动之和。

每个模式都非常简单——只随时间作正弦运动，这一事实使该原理显得很重要。诚然，即使弦的一般运动实际上并不非常复杂，但是存在着另一些系统，例如飞机机翼的抖动，其运动就复杂得多。然而，即使是机翼，我们也发现存在着具有一种频率的某个特定的扭曲方式，以及具有别的频率的其他扭曲方式。如果能够找到这些模式，那么机翼的复杂运动总是可以分解为一些简谐振动的叠加（除非抖动达到这种程度，以至于系统不能再看成是线性的）。

\section{二维波模}

下一个要讨论的例子是二维模式的有趣情况。到目前为止，我们只讲了有关一维的情况——张紧的弦或管中的声波。我们最终应讨论三维情况，但是比较容易的一步是对
\begin{wrapfigure}{l}{0.35\textwidth}
    \centering
    \includegraphics[width=0.33\textwidth]{Chapter49/振动的矩形板}
    \caption{振动的矩形板}
    \label{figure:振动的矩形板}
\end{wrapfigure}
二维情况进行讨论。为了明确起见，我们来讨论一个受到这样约束的矩形橡皮鼓面，它的矩形边界上任何地方的位移都为零，同时设矩形的尺寸为$a$和$b$，如图\ref{figure:振动的矩形板}所示。现在的问题是，可能具有的运动方式有什么特征？我们可以从处理弦问题时所采用的同样步骤着手。如果鼓面根本不受约束，那么可以预期，波将以某种波动形式向前传播。例如，$e^{i\omega t}(e^{-ik_xx+ik_yy})$就代表沿某方向传播的正弦波，此方向取决于$k_x$和$k_y$的相对值。现在，我们怎样才能使$x$轴，即$y=0$的直线成为波节呢？利用从一维弦发展而来的概念，我们可以想象用复函数$(-e^{i\omega t})(e^{-ik_xx-ik_yy})$表示的另一个波。不论$x$和$t$为何值，这两个波的叠加，在$y=0$处得到的位移将为零（虽然这些函数是对负的$y$定义的，在那里没有鼓面在振动，但是这一点可以不予理会，因为在$y=0$处位移确实等于零）。在这种情况下，我们可以视第二个函数为反射波。

然而，我们不仅要求$y=b$处为波节线，而且还要求$y=0$处为波节线。怎么做到这一点呢？问题的解与我们研究晶体反射时所得出的结论有关。只有当$2b\sin\theta$是$\lambda$的整数倍时，这些在$y=0$处相互抵消的波才会在$y=b$处也相互抵消，式中$\theta$是图\ref{figure:振动的矩形板}中所示的角
\begin{equation}
    \label{Eq:I:49:7}
    m\lambda = 2b\sin\theta ~~~~ (m = 0,~1,~2,~\cdots)
\end{equation}

现在，按同样的方法，加上两个函数$-e^{i\omega t}(e^{+ik_xx+ik_yy})$和$+e^{i\omega t}(e^{+ik_xx-ik_yy})$，分别代表来自直线$x=0$的另两个波中的反射波，我们就可以使$y$轴成为波节线。使$x=a$成为波节线的条件与使$y=b$成为波节线的条件相似。即$2a\cos\theta$也必须是$\lambda$的整数倍
\begin{equation}
    \label{Eq:I:49:8}
    n\lambda = 2a\cos\theta
\end{equation}
于是最终的结果是框中来回反射的波展现出一个驻波图样，即一个确定的模式。

因此，要得到一个模式，就必须满足上述两个条件。我们先来求一下波长。由式（\ref{Eq:I:49:7}）和（\ref{Eq:I:49:8}）消去角度$\theta$，就可以求得用$a$，$b$，$n$和$m$表示的波长。最容易的计算方法是将上述两个方程的两边分别用$2b$和$2a$去除，然后将它们平方，再把这两个方程相加。结果是$\sin^2\theta +\cos^2\theta=1=(n\lambda /2a)^2 +(m\lambda/2b)^2$，由此可以解出$\lambda$
\begin{equation}
    \label{Eq:I:49:9}
    \frac{1}{\lambda^2} = \frac{n^2}{4a^2} + \frac{m^2}{4b^2}
\end{equation}
这样，我们根据两个整数决定了波长，同时由波长，可立刻得到频率$\omega$，因为我们知道，频率等于$2\pi c$ 除以波长。

这个结果很有趣并且相当重要，因此我们还应通过纯粹数学的分析，而不是仅从关于反射的讨论来推导出这一结果。让我们用这样选定的四个波的叠加来表示振动，使四条直线$x=0$，$x=a$，$y=0$和$y=b$全都是波节。另外，我们还要求所有的波具有同样的频率，以使合成的运动代表一个模式。由前面对光反射的讨论，我们知道$e^{i\omega t}(e^{-ik_xx+ik_yy})$代表沿图\ref{figure:振动的矩形板}所示的方向传播的波。式（\ref{Eq:I:49:6}），即$k =\omega /c$仍然成立，只要
\begin{equation}
    \label{Eq:I:49:10}
    k^2 = k_x^2 + k_y^2
\end{equation}
显然，由图可知$k_x =k\cos\theta$，$k_y =k\sin\theta$。

现在，我们的矩形鼓面的位移（比如说$\phi$）的方程具有这种很长的形式

\begin{equation}
    \label{Eq:I:49:11a}
    \tag{49.11a}
    \phi = e^{i\omega t} \left[e^{(-ik_xx+ik_yy)}-e^{(+ik_xx+ik_yy)}-e^{(-ik_xx-ik_yy)}-e^{(+ik_xx-ik_yy)}\right]
\end{equation}

虽然这个式子看起来相当混乱，但是计算上述各项之和并不十分困难。将各指数项合并，就可以得到正弦函数，这样一来，位移就是
\begin{equation}
    \label{Eq:I:49:11b}
    \tag{49.11b}
    \phi = \left[4\sin k_xx \sin k_yy\right]\left[e^{i\omega t}\right]
\end{equation}
换句话说，这是一种正弦振动，其图样在$x$方向和$y$方向上也都是正弦的，一切都很好。当然，我们的边界条件在$x=0$和$y=0$处是满足的。我们还要求$\phi$在$x=a$和$y=b$时也等于零。因此，我们必须加进另外两个条件：$k_x a$必须是$\pi$的整数倍，而$k_y b$必须是$\pi$的另一个整数倍。因为我们已经看到$k_x =k \cos\theta$，$k_y =k \sin\theta$，所以立即就得到式（\ref{Eq:I:49:7}）和（\ref{Eq:I:49:8}），并由这两个式子得出最后结果式（\ref{Eq:I:49:9}）。

\setcounter{equation}{11}

现在我们取一个宽是高的两倍的矩形作为例子。如果我们选取$a=2b$，并利用式（\ref{Eq:I:49:4}）和（\ref{Eq:I:49:9}），那么就能够计算出所有模式的频率
\begin{equation}
    \label{Eq:I:49:12}
    \omega^2 = \left(\frac{\pi c}{b}\right)^2 \frac{4m^2+n^2}{4}
\end{equation}
表\ref{Table:49:1}列出几个简单的模式，并且还定性地表示了它们的形状。

\begin{table}[H]
    \caption{}
    \label{Table:49:1}
    \includegraphics[width=0.95\textwidth]{Chapter49/Table/49-1}
\end{table}

关于这种特殊情况，要强调的最重要之点是频率互相不成倍数，它们也不是任何数的倍数。固有频率间具有谐和地相关的概念并不是普遍正确的。对于一维以上的系统它是不正确的，对于比具有均匀密度和均匀张力的弦更复杂的一维系统，它也是不正确的。后一种情况的一个简单例子是悬挂着的链条，链条顶部的张力比底部的张力大一些。如果使这种链条作谐振动，那么就有各种不同的模式和频率，而各频率都不是任何数的简单倍数，模式形状也不是正弦的。

更复杂一些的系统的波模就更加复杂了。例如，在我们的口腔内，声带上部是腔体，依靠使舌头和嘴唇活动，等等，可以做成不同口径和形状的开管或闭管。这是一个极其复杂的共振器，但它不过是一个共振器。当人们运用声带讲话时，声带就用来产生某种类型的音调。这种音调是相当复杂的，出来的有许多声音，但是由于口腔具有各种各样的共振频率，它进一步改变了这个音调。例如，一个歌手能够在同样的音调上发出各种元音a、o或oo等，但是它们听起来不一样，因为在口腔中，各种谐音的共鸣程度不同。空腔的共振频率在改变发声中的极端重要性可以用一个简单的实验来演示。由于声音的传播速率随介质密度的平方根的倒数而变化，因此利用不同的气体可以改变声音的速率。如果我们用氦气代替空气，这样密度就降低，声音的传播速率大大增加，空腔的所有频率都将升高。因而，如果在某人讲话前用氦气充满肺部，那么虽然声带还是以同样的频率振动，但他的发声特征将发生极大变化。

\section{耦合摆}

最后，我们要着重指出，不仅仅是复杂的连续系统存在着波模，而且非常简单的力学
\begin{wrapfigure}{r}{0.35\textwidth}
    \centering
    \includegraphics[width=0.33\textwidth]{Chapter49/两个耦合摆}
    \caption{两个耦合摆}
    \label{figure:两个耦合摆}
\end{wrapfigure}
系统也存在着波模。前一章中讨论过的两个耦合摆组成的系统，就是一个很好的例子。在那一章中已经证明，可以把这种运动看作具有不同频率的两个谐运动的叠加来分析。因此，即使是这样的系统也可以用谐运动或波模来分析。弦具有无穷个模式，二维的表面也具有无穷个模式。如果我们知道如何计算无穷大，那么在某种意义上，二维表面的模式数目是一个二重无穷大。但是，仅有两个自由度，只要求用两个变量来描述的一个简单的力学系统只有两个模式。

让我们在摆是等长的情况下对这两个模式进行数学的分析。如图\ref{figure:两个耦合摆}所示，令一个摆的位移为$x$，另一个摆的位移为$y$。如果没有弹簧，由于重力，作用在第一个物体上的力与该物体的位移成正比。如果没有弹簧，仅仅对于这一个摆来说，就应该有某个固有频率$\omega_0$。没有弹簧时，运动方程应为
\begin{equation}
    \label{Eq:I:49:13}
    m\frac{d^2x}{dt^2} = -m\omega_0^2 x
\end{equation}

如果没有弹簧，另一个摆应以同样的方式摆动。但当存在弹簧时，除去由于重力而产生的恢复力外，还有一个附加力在拉第一个物体。这个拉力取决于$x$超过$y$的距离，并且与其差额成正比，因此，拉力就等于取决于几何位置的某个常数乘以$(x-y)$。同样大小但方向相反的力作用在第二个物体上。因此，要解的运动方程是
\begin{equation}
    \label{Eq:I:49:14}
    m\frac{d^2x}{dt^2} = -m\omega_0^2 x - k(x-y),~~ m\frac{d^2y}{dt^2} = -m\omega_0^2 y - k(y-x)
\end{equation}

为了求出两个物体都以同样频率摆动的运动，我们必须决定每一个物体移动多少。换句话说，摆$x$和摆$y$将以同样的频率振动，但是它们的振幅必须具有确定的数值$A$和$B$，$A$和$B$的关系是固定的。我们用这种解来尝试一下
\begin{equation}
    \label{Eq:I:49:15}
    x=Ae^{i\omega t},~~y=Be^{i\omega t}
\end{equation}
如果将这两个试解代入方程式（\ref{Eq:I:49:14}），合并同类项，结果为
\begin{equation}
    \label{Eq:I:49:16}
    \begin{aligned}
        \left(\omega^2 - \omega_0^2 - \frac{k}{m}\right)A = -\frac{k}{m}B \\
        \left(\omega^2 - \omega_0^2 - \frac{k}{m}\right)B = -\frac{k}{m}A 
    \end{aligned}
\end{equation}
所写出的方程中，公共因子$e^{i\omega t}$已被消去，并且都除以$m$。

现在，我们看到，有两个似乎包含两个未知数的方程式。但是事实上没有两个 未知数，因为运动的整个大小是无法由这些方程来决定的。上述方程可以决定的仅仅是$A$与$B$之比 ，但是这两个方程必定给出同样的比例。 要使这两个方程相一致，就要求频率是某个非常特殊的值。

在这种特殊情况下，频率可以相当容易地求出。如果将这两个方程相乘，那么结果是
\begin{equation}
    \label{Eq:I:49:17}
    \left(\omega^2 - \omega_0^2 - \frac{k}{m}\right)^2AB = \left(\frac{k}{m}\right)^2AB 
\end{equation}
方程两边的$AB$项可以消去，除非$A$和$B$都为零，而那就意味着根本没有运动。如果存在运动，那么剩下的两项必定相等，这给出一个要求出解的二次方程。结果是存在两个可能的频率
\begin{equation}
    \label{Eq:I:49:18}
    \omega_1^2 = \omega_0^2,~~\omega_2^2 = \omega_0^2 + \frac{2k}{m}
\end{equation}
另外，如果把这些频率值代回到方程（\ref{Eq:I:49:16}）中，我们发现对第一个频率有$A=B$，对第二个频率有$A=-B$。这些就是“模式形状”，这很容易用实验来验证。

显然，在$A=B$的第一种模式中，弹簧从来没有被拉伸，因而两个摆都以频率$\omega_0$振动，就好像弹簧不存在时一样。在$A=-B$的另一个解中，弹簧提供一个恢复力，因而提高了振动频率。如果两个摆具有不同的长度，则可得出更为有趣的情况。有关的分析与上面给出的十分相似，我们将它留给读者作为练习。

\section{线性系统}

现在，我们来总结一下上面所讨论过的概念，它可能包含了数学物理的最普遍、最奇妙的原理的所有方面。如果我们有一个特征与时间无关的线性系统，那么它的运动并不一定具有什么特别的简单性，事实上，它可以是非常复杂的，但是存在着一些（通常是一系列）非常特殊的、其整个图样随时间指数地变化的运动。对我们现在正在谈论的振动系统，指数是虚数，因而我们不说它随时间“指数地”变化，而宁可说它随时间“按正弦规律”变化。然而，人们可以更一般地说运动随时间以非常特殊的形状及非常特殊的模式指数地变化。系统最普遍的运动总是可以表示成包含各个不同指数的运动的叠加。

对于正弦运动的情况，将上述结论再叙述一遍是值得的：一个线性系统不一定作纯粹的正弦运动，即以确定的单个频率运动，但是无论它如何运动，这个运动总可以表示成纯正弦运动的叠加，这些正弦运动中每个运动的频率是系统的一个特征，并且每个运动的图样或波形也是系统的一个特征。在任何一个这样的系统中，一般运动的特征可以这样来表示，即给出各个模式的强度和相位，并把它们统统合在一起。叙述这一点的另一种方式是：任何一个线性振动系统等效于一组互相独立的谐振子，每个谐振子具有与各模式相应的固有频率。

我们评述一下模式与量子力学的联系作为这一章的结束。量子力学中的振动客体（或在空间变化的东西）是概率函数的振幅，它给出了在一个已知组态中，找到一个电子或一组电子的概率。这个振幅函数可以随空间和时间变化，并且事实上还满足一个线性方程。然而，在量子力学中有一种变换，在这种变换中我们所谓的概率振幅的频率就等于经典概念中的能量。因此，只要用能量 一词来代替频率 ，我们就可以将上述原理变成适用于量子力学的情况。它变成这样：一个量子力学系统，例如一个原子，不必具有一个确定的能量，正像一个简单的力学系统并不一定要具有一个确定的频率一样。但是，无论系统表现得怎样，其行为总可以表示为一些具有确定能量的状态的叠加。每个状态的能量是该原子的一个特征，同时决定在不同地点找到粒子概率的振幅图样也是原子的一个特征。一般的运动可以用给出各个不同能量状态的振幅来描写。这就是量子力学中能级的由来。由于量子力学是用波来表示的，在电子不具备足够的能量最终脱离开质子的情形下，这些波都是约束波 。就像弦的约束波一样。量子力学波动方程的解中存在许多确定的频率。量子力学解释说这些就是确定的能量 。因此一个量子力学系统，由于它是用波来表示的，可以具有能量固定的确定状态；各种各样原子的许多能级就是例子。